% Wampanoag Multi-Source — Source Workflow
% Issue #1365
% Status: Research complete, implementation pending
% !TEX program = xelatex
% !TEX encoding = UTF-8 Unicode

\section{Wampanoag: Multi-Source (Wood, Winslow, Shepard, Eliot)}
\label{sec:1365}

\subsection{Source Description}

This section covers all historical Wampanoag/Massachusett source vocabularies
beyond the Trumbull compilation (covered in \S\ref{sec:1363}). Four primary
sources span the critical transition from raw English-hearer perception through
systematic missionary orthography:

\begin{description}
  \item[William Wood (1634)] \emph{New Englands Prospect}, \textasciitilde262
    words. Pre-Eliot English ad-hoc orthography. Wood was an English colonist
    living in the Massachusetts Bay Colony with no linguistic training. His
    transcriptions reflect the earliest English-hearer perception of Wampanoag
    sounds, before any regularization system existed.
  \item[Edward Winslow (1624)] \emph{Good Newes from New England}, 22 scattered
    words and phrases. The earliest published Wampanoag vocabulary. Winslow was
    a Mayflower passenger and Plymouth Colony governor who described the
    language as ``very copious, large, and difficult'' and was unusually honest
    about his linguistic limitations.
  \item[Thomas Shepard (1647)] \emph{The Day-Breaking, if not the sun-rising
    of the Gospell with the Indians in New-England}, 11 native words plus two
    one-sentence prayers. Predates Eliot's grammar by 19 years and his Indian
    Primer by 22 years. The prayers represent Eliot's earliest work in the
    language, before he had fully analyzed its grammar.
  \item[John Eliot (1666)] \emph{The Indian Grammar Begun}, the first published
    grammatical analysis of any Algonquian language. Eliot's orthography is a
    deliberately constructed system, not a casual English-based spelling. He
    distinguishes aspirated vs.\ unaspirated stops using English voiced/voiceless
    pairs (p\textasciitilde b, t\textasciitilde d, k\textasciitilde g).
\end{description}

\subsection{Orthography Systems}

\subsubsection{William Wood (1634)}

Wood's vocabulary is the largest pre-Eliot Massachusett word list. Its
orthography follows English sound-spelling conventions of the 1630s:

\begin{itemize}
  \item \textbf{Initial /w/ retained:} \emph{Wigwam}, \emph{Web}, \emph{Wompey}
  \item \textbf{$\langle$ea$\rangle$ for /iː/:} \emph{Kean} (cf.\ Eliot's \emph{keen})
  \item \textbf{$\langle$s$\rangle$ for /ʃ/ inconsistently:} \emph{Sawup} where
    later sources might write \emph{sh-}
  \item \textbf{$\langle$k$\rangle$ preferred over $\langle$c$\rangle$:}
    Initial /k/ written \emph{K-} even before front vowels (\emph{Kean},
    \emph{Ksitta})
  \item \textbf{No diacritics:} Wood uses plain English letters only
  \item \textbf{Vowel length not marked:} Short vs.\ long vowel distinction
    invisible
  \item \textbf{Pidginization noted:} Goddard (1977) observes a ``good deal of
    pidginization'' in Wood's data (Salvucci 2002)
\end{itemize}

\subsubsection{Edward Winslow (1624)}

Winslow's orthography is purely English-hearer perception with zero
regularization. He was recording an entirely unfamiliar language with no
reference system. The vocabulary overwhelmingly represents the Pokanoket
dialect of Wampanoag (Massasoit's people), with some possible influence from
the Patuxet band through Tisquantum.

Key characteristics:
\begin{itemize}
  \item \textbf{Pre-Eliot:} No systematic orthography; words are ad-hoc English
    approximations
  \item \textbf{$\langle$h$\rangle$ written reliably:} Winslow writes /h/ more
    consistently than any other early Wampanoag source except Edwards (1787
    Mahican). This is one of the most valuable aspects of his small corpus.
  \item \textbf{$\langle$l$\rangle$ $\rightarrow$ $\langle$n$\rangle$:}
    \emph{Winsnow} for ``Winslow'' --- no /l/ in Wampanoag
  \item \textbf{Pidginization:} Goddard (1977) identifies clear evidence of
    pidginization, including use of transitive inanimate \emph{namen} where
    transitive animate \emph{nunaw} would be grammatically correct
\end{itemize}

\subsubsection{Thomas Shepard (1647)}

The Wampanoag words in this pamphlet show unstandardized English orthography
with the following characteristics:

\begin{itemize}
  \item \textbf{$\langle$ch$\rangle$ for /č/:} \emph{Chechainuppan},
    \emph{Chepian}, \emph{Cheehesom}
  \item \textbf{$\langle$sh$\rangle$ and $\langle$s$\rangle$ for /š/:}
    \emph{Evangenesh}, \emph{Iehovah} (given English biblical name spelling)
  \item \textbf{$\langle$w$\rangle$ for /w/:} \emph{Wigwam}, \emph{Wowein},
    \emph{wicke}
  \item \textbf{Final -an for animate nouns:} \emph{Chepian},
    \emph{Chechainuppan}
  \item \textbf{$\langle$oo$\rangle$ for /uː/:} \emph{Noonatomen}
  \item \textbf{$\langle$p$\rangle$ for /b/:} \emph{Chepian} (cf.\ Eliot's
    \emph{chepiân})
  \item \textbf{Preaspiration:} \emph{cke} in \emph{wicke} may indicate
    preaspirated stop
\end{itemize}

The pamphlet contains the earliest recorded Wampanoag Christian prayers:

\begin{quote}
  \emph{Amanaomen Iehovah tahassen metagh} --- ``Take away Lord my Stony heart''
\end{quote}

\begin{quote}
  \emph{Cheehesom Iehovah kekowhogkew} --- ``Wash Lord my soule''
\end{quote}

\subsubsection{John Eliot (1666)}

Eliot's grammar presents a deliberately constructed orthography. His key
innovation is distinguishing aspirated vs.\ unaspirated stops using English
voiced/voiceless pairs (p\textasciitilde b, t\textasciitilde d,
k\textasciitilde g). This is the most important orthographic feature for
phonological reconstruction.

Eliot describes 5 primary vowels with length distinctions:
\begin{itemize}
  \item \textbf{a} (short) vs.\ \textbf{â} (long)
  \item \textbf{e} /iː/ (always long)
  \item \textbf{i} (short) vs.\ \textbf{ee}/\textbf{ie} (long --- inconsistently
    marked)
  \item \textbf{o} (short) vs.\ \textbf{ô} (long)
  \item \textbf{u} (short schwa) vs.\ \textbf{oo} (long /uː/)
\end{itemize}

\subsection{Vowel and Consonant Tables}

\begin{table}[htbp]
\centering
\caption{Comparative Orthography Table}
\label{tab:1365-compare}
\begin{tabular}{llllll}
\hline
English & Wood (1634) & Shepard (1647) & Eliot (1666+) & Cotton (1708) & Trumbull (1903) \\
\hline
man & --- & --- & wosketomp & wosketomp & wosketomp \\
woman & --- & --- & mittamwossis & mittamwo˘ssis & mittámwossis \\
house & Wigwam & Wigwam & wetu & wetu & wétu \\
devil & --- & Chepian & chepiân & chepiân & chepeân \\
king/sachem & --- & Sachem & sôchim & sôchim & sôchim \\
I & Kean & --- & keen & keen & kēn \\
very angry & --- & musquantum & --- & --- & --- \\
good & Winnet & --- & wunne & wunne & wunne \\
white & Wompey & --- & wompi & wompi & wômpei \\
sleep (classifier) & Sawup & --- & --- & --- & --- \\
\hline
\end{tabular}
\end{table}

\begin{table}[htbp]
\centering
\caption{Vowel Representations by Source}
\label{tab:1365-vowels}
\begin{tabular}{llllll}
\hline
Phoneme & Wood & Shepard & Eliot & Cotton & Trumbull \\
\hline
/iː/ & ea, ee & e, ee & e, ee & e, ee & ē \\
/a/ & a & a & a & a & a \\
/aː/ & a, aw & a & â & â & â \\
/uː/ & oo & o, oo & o, ô, ꝏ & o, ô & ô, u \\
/ə/ & --- & a, e & u & u, u˘ & u \\
/o/ & o & o & o & o & o \\
\hline
\end{tabular}
\end{table}

\begin{table}[htbp]
\centering
\caption{Consonant Representations by Source}
\label{tab:1365-consonants}
\begin{tabular}{llllll}
\hline
Phoneme & Wood & Shepard & Eliot & Cotton & Trumbull \\
\hline
/p/ & p & p & p\textasciitilde b & p\textasciitilde b & p \\
/pʰ/ & --- & --- & p & p & p (with dot) \\
/t/ & t & t & t\textasciitilde d & t\textasciitilde d & t \\
/tʰ/ & --- & --- & t & t & t (with dot) \\
/k/ & k & k & k\textasciitilde g & k\textasciitilde g & k \\
/kʰ/ & --- & --- & k & k & k (with dot) \\
/č/ & ch & ch & ch & ch & tc \\
/š/ & sh & sh & sh & sh & sh \\
/w/ & w & w & w & w & w \\
/j/ & --- & --- & y & y & y \\
\hline
\end{tabular}
\end{table}

\subsection{Source-Specific Orthographic Patterns}

Each source has characteristic ``errors'' or systematic deviations from a
standardized phonemic representation:

\begin{description}
  \item[Wood (1634)] No vowel length; English digraphs predominate; treats
    unaspirated stops as voiced; pidginized forms; word-initial \emph{a-}
    prefix frequently attached to roots.
  \item[Shepard (1647)] No consistent system; words are ad-hoc English
    approximations; preaspiration inconsistently shown; $\langle$ck$\rangle$
    for /k/.
  \item[Eliot (1666+)] Most consistent system; treats unaspirated stops as
    voiced English letters (b, d, g) but means unaspirated; vowel length
    inconsistently marked on some vowels; macron on long vowels (â, ô);
    distinctive $\langle$ꝏ$\rangle$ (oo ligature) for /uː/.
  \item[Cotton (1708)] Extends Eliot's system; adds breve for ultra-short
    vowels; shows late shift in some vowels; more consistent vowel marking.
  \item[Trumbull (1903)] Standardized historical reconstruction; marks length
    and aspiration consistently; normalizes variation.
\end{description}

\subsection{Problems Encountered}

\begin{itemize}
  \item \textbf{Source provenance correction:} The ``Anonymous 1647'' vocabulary
    in ALR Vol.\ 27 is actually Thomas Shepard's \emph{The Day-Breaking}.
    Furthermore, the sample words in the original ALR listing (\emph{kean},
    \emph{kenomptom}, \emph{wean}, \emph{nukus}, \emph{kodtomp}, \emph{ompmom},
    \emph{sawop}, \emph{acheta co}) do NOT match the words actually found in
    the 1647 text. These words originate from William Wood's 1634 vocabulary.
    The correct Shepard word list contains 11 items including the two prayers.
  \item \textbf{Pre-Eliot vs.\ post-Eliot:} Sources span the transition from
    raw English-hearer perception (Wood, Winslow) through systematic missionary
    orthography (Eliot). The same letter sequence may have different phonetic
    values in different sources.
  \item \textbf{Eliot's stop aspiration system:} Eliot is the only colonial
    source who systematically distinguishes aspirated from unaspirated stops
    using orthographic pairs (b/p, d/t, k/g). His $\langle$b, d, g$\rangle$
    represent unaspirated lenis stops (not voiced), while $\langle$p, t,
    k$\rangle$ represent aspirated fortis stops.
  \item \textbf{Pidginization:} Goddard (1977) identifies pidginization in both
    Wood and Winslow. If the English were learning a simplified contact variety,
    word order may not reflect full Wampanoag syntax, verb forms may be bare
    stems rather than inflected forms, and vocabulary may be mixed from multiple
    Algonquian languages.
  \item \textbf{Intermediary filter:} Winslow's data passed through Tisquantum
    (Squanto) and Hobbamock as interpreters. Tisquantum's English was learned
    in London, not from the Pilgrims, and his Wampanoag was Patuxet, a
    sub-dialect closely related to but not identical with Pokanoket.
  \item \textbf{Dialect uncertainty:} Pokanoket (Winslow/Wood) vs.\ Natick
    (Eliot/Trumbull) differences are unresolved. The distinction may represent
    real dialect differences, register differences (pidgin vs.\ full language),
    chronological differences (1620s vs.\ 1660s), or orthographic convergence.
  \item \textbf{Vowel instability:} Wood's triple spelling
    \emph{abbomocho}/\emph{abomocho}/\emph{abamacho} for ``devil'' is the
    strongest evidence in the entire SNEA corpus for colonial vowel instability.
  \item \textbf{Preaspiration /h/ unreported:} Like all pre-Eliot sources, Wood
    does not write preaspirated /h/ before stops. Winslow is the exception ---
    he writes /h/ more consistently than any other early source.
\end{itemize}

\subsection{Research Approach}

The multi-source parsing strategy follows a four-stage pipeline:

\begin{enumerate}
  \item \textbf{Source detection:} Identify which source's orthographic
    fingerprint is present. Each source has a characteristic orthographic
    texture that should be identifiable from spelling patterns.
  \item \textbf{Source normalizer:} Convert to Eliot-consistent form. Apply
    source-specific decoding rules --- the same letter sequence may have
    different phonetic values in different sources.
  \item \textbf{Eliot parser:} Apply Grammar rules to produce a phonemic
    representation. Eliot's orthography is the most systematic and serves as
    the intermediate representation.
  \item \textbf{Unified output:} Produce standard modern transcription.
\end{enumerate}

\subsubsection{Source-Adaptive Parsing Strategy}

\begin{description}
  \item[Wood] English sound-spelling $\rightarrow$ phoneme via lookup table.
    Key mappings: ea$\rightarrow$/iː/, oo$\rightarrow$/uː/, k$\rightarrow$/k/,
    s$\rightarrow$/s\textasciitilde š/.
  \item[Shepard] Ad-hoc English $\rightarrow$ phoneme via context. Same as Wood
    but with higher ambiguity.
  \item[Eliot] Eliot orthography $\rightarrow$ phoneme via Grammar rules.
    p/b$\rightarrow$/p/ vs.\ /pʰ/, t/d$\rightarrow$/t/ vs.\ /tʰ/,
    k/g$\rightarrow$/k/ vs.\ /kʰ/, â$\rightarrow$/aː/, ô$\rightarrow$/oː\textasciitilde uː/.
  \item[Cotton] Same as Eliot + breve $\rightarrow$ extra-short vowel. Vowel
    quantity more reliable.
  \item[Trumbull] Modern normalized $\rightarrow$ phoneme directly. The cleanest
    mapping.
\end{description}

\subsubsection{Resolving Ambiguity}

Where a spelling could represent multiple phonemes (e.g., Wood's $\langle$ea$\rangle$
= /iː/ or /ɛː/), the parser uses source context and Algonquian comparative
evidence. The Proto-Algonquian reconstructions in Aubin (1975) provide the
external check needed to disambiguate competing analyses.

\subsection{Conversion Workflow}
\texttt{[Pending implementation --- see Issue \#1365]}

The conversion pipeline will be implemented as a Python module in
\texttt{src/commons/parsing/}. The module will:

\begin{itemize}
  \item Detect source fingerprint from orthographic patterns
  \item Apply source-specific normalizer functions
  \item Route through the unified Eliot parser
  \item Output standard modern transcription
  \item Handle the pre-Eliot vs.\ post-Eliot transition
  \item Resolve ambiguity using PA comparative evidence
\end{itemize}

\subsection{Linguist Feedback Template}

\texttt{[Template --- content pending review]}

The following questions are open for linguist review:

\begin{itemize}
  \item Are the source-specific vowel and consonant mappings correct?
  \item Does the proposed Eliot-consistent intermediate representation capture
    the phonemic distinctions correctly?
  \item Are there additional sources that should be included?
  \item Is the pidginization analysis (Goddard 1977) accepted, or are there
    alternative interpretations?
  \item How should the Pokanoket/Natick dialect distinction be handled in the
    unified output?
\end{itemize}

\subsection{Related Work: Colonial Recorder Perception Problem}

The core challenge of this research card --- English-speaking recorders
misperceiving and mis-transcribing Algonquian phonemes --- is a known
phenomenon documented across Eastern North America. Two researchers' work is
directly relevant:

\textbf{Dr.\ Keith Cunningham} (Linguist, Nanticoke Indian Tribe). His doctoral
dissertation \emph{A Phonological Analysis of Nanticoke With Practical
Applications for Language Revitalization} (Georgetown University) analyzes three
18th-century Nanticoke word lists (Heckewelder 1785) using the same methodology:
historical phonology from colonial recorders, with the same challenges of
multilingual informants and orthographic complexity. His 2024 Algonquian
Conference presentation \emph{A Revised Phonological Analysis of the Heckewelder
Vocabulary of Nanticoke} directly parallels the multi-source Wampanoag analysis
--- both involve multiple colonial recorders (Wood, Winslow, Shepard, Eliot for
Wampanoag; Heckewelder, Du Ponceau for Nanticoke) with different orthographic
conventions that must be cross-calibrated. Cunningham's work demonstrates that
the colonial recorder perception problem extends beyond SNEA into the broader
Eastern Algonquian family.

\textbf{Dr.\ Craig Kopris} (Independent scholar). His paper \emph{Les sons
wyandots perçus par des oreilles étrangères} (Wyandot sounds perceived by
foreign ears, \emph{Recherches amérindiennes au Québec}, 2014) is the exact
framing of the multi-source Wampanoag problem: how European recorders with
different linguistic backgrounds (17th c.\ English, 18th c.\ English, 19th c.\
German-influenced American English) systematically misperceived and
mis-transcribed indigenous phonemes. His \emph{Wyandot Phonology: Recovering the
Sound System of an Extinct Language} applies the same recovery methodology to an
unrelated language family (Iroquoian), demonstrating that the foreign-ear
perception problem is universal across colonial documentation of North American
languages.

The multi-source Wampanoag corpus (Wood 1634, Winslow 1624, Shepard 1647,
Eliot 1666) spans the full range of colonial recorder competence --- from raw
English-hearer perception (Wood, Winslow) through systematic missionary
orthography (Eliot). This range makes it the best case study for the recorder
calibration problem that both Cunningham and Kopris address in their respective
language families.

\subsection{Key References}

\begin{itemize}
  \item Aubin, George F. 1975. \emph{A Proto-Algonquian Dictionary}. Ottawa:
    National Museums of Canada.
  \item Aubin, George F. 1978. ``Toward the Linguistic History of an Algonquian
    Dialect: Observations on the Vocabulary.'' In William Cowan (ed.),
    \emph{Ninth Algonquian Conference}, pp.\ 1--9.
  \item Clark, Michael P. (ed.). 2003. \emph{The Eliot Tracts}. Westport, CT:
    Praeger.
  \item Cogley, Richard W. 1999. \emph{John Eliot's Mission to the Indians
    before King Philip's War}. Cambridge, MA: Harvard University Press.
  \item Cotton, Josiah. 1829. \emph{Vocabulary of the Massachusetts (or Natick)
    Indian Language}. Cambridge: Charles Folsom.
  \item Eliot, John. 1666. \emph{The Indian Grammar Begun}. Cambridge:
    Marmaduke Johnson.
  \item Goddard, Ives. 1977. ``The Pidginization of Massachusetts: Evidence
    from Wood's Vocabulary.'' \emph{International Journal of American
    Linguistics} 43: 222--226.
  \item Goddard, Ives \& Bragdon, Kathleen J. 1988. \emph{Native Writings in
    Massachusett}. Philadelphia: American Philosophical Society.
  \item Pilling, James C. 1891. \emph{Bibliography of the Algonquian Languages}.
    Washington, D.C.: Government Printing Office.
  \item Salvucci, Claudio (ed.). 2002. \emph{ALR Volume 27: Wood's Vocabulary of
    Massachusett}. Bristol, PA: Evolution Publishing.
  \item Shepard, Thomas (attrib.). 1647. \emph{The Day-breaking, if not the
    sun-rising of the Gospell with the Indians in New-England}. London: Rich.
    Cotes.
  \item Trumbull, J. Hammond. 1903. \emph{Natick Dictionary}. Washington, D.C.:
    Bureau of American Ethnology.
  \item Williams, Roger. 1643. \emph{A Key into the Language of America}.
    London: Gregory Dexter.
\end{itemize}
